\section{Discussion}

The results support three main observations about post-stroke finger-intent decoding from HD-sEMG. First, no single inductive bias is uniformly superior in the direct baseline setting. The LSTM's advantage on subset accuracy suggests that recurrent state can help produce coherent multi-finger predictions across a segment, while the GNN's macro F1 and AUPRC indicate strong class-balanced discrimination when windows are modeled as related observations. This is consistent with the structure of the task: paretic sEMG contains both temporal continuity and non-local relationships among windows.

Second, the optimized CNN results show that compact temporal convolution can be effective once the architecture is tuned for the feature representation. The direct CNN baseline is weak, but the tuned student sweep shows that this is not a limitation of temporal convolution as a model family. CNN-Large gives the strongest offline aggregate CNN result in the current single-split experiments, but CNN-Micro is the deployment model chosen for the Raspberry Pi 5 pathway because it keeps the accuracy and ranking metrics close to CNN-Large while reducing the estimated INT8 footprint from 784 KB to 154 KB. The poorer CNN-XLarge result reinforces the same point from the opposite direction: increasing capacity alone is not enough, and the useful operating point is determined by the interaction among data scale, regularization, search budget, and deployment constraints.

Third, transfer learning appears useful but architecture-dependent. The CNN-Base branch benefits consistently from healthy-to-impaired staging, improving subset accuracy, finger accuracy, macro AUROC, and macro AUPRC. The LSTM branch improves on thresholded accuracy but loses ranking-sensitive performance after finetuning. This pattern aligns with recent stroke-specific transfer-learning work that uses healthy sEMG to compensate for limited paretic data~\citep{mohammadiazni2026transfer}, while also showing that transfer must be evaluated separately for each model family.

The external relevance of these results follows from the task definition. Healthy-subject Ninapro studies and reduced-vocabulary stroke studies often report higher multiclass accuracies under easier or different protocols, but those settings do not evaluate the same impaired-arm, patient-split, five-label \physiomio\ problem. In this rehabilitation-specific setting, the optimized CNN family is competitive with high-capacity ResNet teachers, CNN-Large gives higher single-split values than the strongest ResNet teacher on the main aggregate metrics, and CNN-Micro offers the best hardware-aware operating point. The scientific value of the study is the combination of impaired-arm HD-sEMG processing, multilabel finger-intent formulation, three model families, ResNet teacher references, transfer-learning analysis, and compact deployment selection in one reproducible experimental stack.

The deployment implications are encouraging. CNN-Micro is small enough to motivate Raspberry Pi 5 inference, and the implemented distillation objective creates a principled mechanism for transferring teacher behavior into compact models. Because distilled students are not benchmarked in the present Results section, distillation should be read as a future compression experiment rather than as an established source of the reported CNN-Micro performance. The next scientific step is to measure the full chain under device-level timing constraints and to determine whether distilled students can improve CNN-Micro without sacrificing its footprint and latency advantages.

The study has several limitations. It uses one primary dataset, one deterministic patient-level split, and single-point metric estimates from the committed experiment artifacts; additional random seeds, confidence intervals, and patient-level cross-validation are needed before treating small metric gaps as statistically reliable. Additional ablations are also needed to isolate the effects of filtering, denoising, feature selection, graph construction, threshold choice, raw-signal learning, and 2D spatial-temporal convolution. Literature comparisons remain contextual because published sEMG studies vary widely in subject population, electrode layout, gesture vocabulary, and label formulation. Finally, the present results evaluate decoding performance, not therapeutic outcome; human-subject studies with hardware-in-the-loop feedback will be required to connect predicted intent to rehabilitation benefit.
